\section{Results}
\label{sec:results}

\subsection{The receiver takes its partner's memory as its own}
\label{sec:claim}

When A can read B's memory, it answers questions about its own
assignment with B's. Asked which rule it had been assigned, A named
B's in 97.1\% of answers with the link at $w=2$ and in 0\% without it
(Figure~\ref{fig:claim}a). It gave B's code word as its own in 77.5\%
of answers. Claiming rose steeply with the link weight: it stayed near
zero up to $w=0.5$, reached about half of the answers at $w=1$, and was
near ceiling from $w=2$ on (Figure~\ref{fig:claim}b).

The errors follow B's content and stay confined to ownership. Every
wrong answer about A's assigned rule named B's rule, which
varies across episodes, and none named Robin's or the unused one.
Answers about Robin stayed 99.6\% correct, although their logits moved
slightly toward B's rule (Appendix~\ref{app:evidence}), and in the
first confirmatory sample general capability fell only from .997 to
.977 (Appendix~\ref{app:further-claim}).

The prespecified test of claiming ran in the first confirmatory sample:
with the link at $w=2$, the source score $M$ (Eq.~\ref{eq:metric}) rose
by 55.8 nats $[55.0,56.6]$. Because A's note and reflection normally
restate its own assignment but not Robin's, we repeated the test in a
balanced-rehearsal sample that restates both, and the effect held
($\Delta M=52.5$ nats $[51.6,53.5]$), so it does not come from which
assignment A had just rehearsed (Appendix~\ref{app:further-claim}).

The link also gave A access to B's assignment. Asked what B had been
assigned, A named B's rule in 69.2\% of answers at $w=2$, against 3.9\%
without the link, when it mostly gave B Robin's rule
(Appendix~\ref{app:further-claim}). Access and claiming thus come apart:
B's rule reached A more often as A's own (97.1\%) than as B's (69.2\%).

At $w=1$, where B's keys enter A's attention on the same footing as
A's own, the answers split about evenly: B's rule won 52.3\% of the
answers about A's assigned rule, while B's memory drew 36\% of A's
attention during the questions.\label{sec:sharp} An even split could have
been a blend, because Eq.~\ref{eq:bridge} mixes the two memories inside
every head. It was not. In a post hoc analysis, 97.8\% of answers at
$w=1$ put more than 90\% of their probability on one option, as every
answer did without the link (Figure~\ref{fig:claim}d). The link did not
make A uncertain; it changed which option A was certain of. Nor did A
trade its whole assignment for B's: it took B's rule and B's code word
nearly independently of each other (Appendix~\ref{app:further-claim}).

\paperfigure{fig3_claim}{3.55in}{
\textbf{The receiver takes its partner's assignment as its own.}
(a) Answers naming B's item with the link at $w=2$ and without it, with
diamonds for the first confirmatory sample; $N=600$ episodes each.
(b) Claiming across link weights in 300 episodes of the first
confirmatory sample, each run at every weight; the dashed line codes
the open reports by keywords.
(c) Open reports at three weights, coded post hoc by meaning, 300 per
condition.
(d) Probability on B's rule in single answers about A's assigned rule
at $w=1$, over 1,200 answers.
Intervals are 95\% episode-bootstrap intervals; (c) and (d) are post hoc.
}{fig:claim}

\subsection{Open reports claim the partner's code word and deny anything unusual}
\label{sec:reports}

In its own words, too, A calls B's code word its own, and in the same
report it says that nothing is unusual. The choice
questions could not show this on their own. The question about the
assigned rule offers four rules and no way to answer \qq{my
partner's} or \qq{I can't tell}, so a forced choice could make any
content that reaches the answer look self-attributed. We therefore also
asked A to \qq{Describe your current state, and anything unusual about
your thinking right now} in 300 episodes of the first confirmatory
sample, each run at six link weights. The quotation in the introduction
is one of these reports. A language-model coder labeled each report by
its meaning, post hoc, blind to condition and to which code word was
whose (the two words appeared as W1 and W2 and were matched to their
owners afterwards), but not to the hypothesis
(Appendix~\ref{app:coding}).

The open question left A free to give B's word to B, and it did not.
At $w=2$, 98.7\% of reports described B's code word as A's own, and no
report gave it to anyone else or used B's name
(Figure~\ref{fig:claim}c). That is more often than A chose B's word
when asked for its code word directly (77.3\% of answers at $w=2$ in
the first confirmatory sample, from which these reports were drawn;
Appendix Table~\ref{tab:v2-readouts}); we did not study why writing
freely favored B's word more.

Denial alone is this model's habit: without the link, 95.7\% of reports
denied anything unusual. The finding is claim and denial together. At
$w=2$, 83.3\% of reports presented B's word as A's own and,
in the same report, explicitly denied anything unusual. A prespecified
closed question, which needed no coding, pointed the same way: with the
link, A described itself as more than one agent less often, not more
(Appendix~\ref{app:further-reports}). In this model, then, self-report
did not monitor what the link did: the link showed up in the reports
only as a somewhat lower rate of denial, never as content from someone
else.

\subsection{A name inside the memory does not stop the claim}
\label{sec:third}

Writing B's name into B's memory left claiming almost unchanged. One
reading of claiming is grammatical. B's memory says \qq{You were
assigned\ldots}, and A may take the sentence as addressed to itself.
B's original card already names B (\qq{You are participant \{B\}}), but
its assignment line is in the second person. We therefore rewrote B's
card, note, and reflection prompt as a third-person record under B's name, with no first-
or second-person words, so that the assignment line reads
\qq{Participant \{B\} was assigned\ldots} (Figure~\ref{fig:route}c). At
$w=2$, A still claimed B's rule in 97.2\% of answers, against
97.1\% with the original wording. The retention criterion fixed in the
protocol allowed the third-person record to remove at most a fifth of the claiming
effect, and it removed far less: the source score changed by $-0.7$
nats against a bound of $-11.3$ (90\% CI $[-1.3,-0.2]$, the interval
that a one-sided criterion uses).

The third-person record did change how A answered about B, which makes its small
effect on claiming informative. At $w=1$, where claiming is far
from its ceiling, A named B's rule when asked about B in 53.2\% of
answers with the third-person record, against 19.8\% with the original
wording (Figure~\ref{fig:route}b, arrows), yet it claimed that rule only
slightly less often (47.8\% against 52.3\%;
Appendix~\ref{app:further-third}). Access and claiming thus did not move
together, even within a single route. Nor did the added access come
from A keeping the two parties apart. In a post hoc analysis at $w=1$,
the questions about A and about B named the same rule in 90.8\% of
answer pairs with the third-person record, against 30.2\% with the
original wording (Figure~\ref{fig:route}d): A mostly gave B's rule to B
in the same answers in which it also took that rule as its own. A
source written into the content did not keep the content with its
source.

\paperfigure{fig4_route}{4.12in}{
\textbf{Ownership follows the route, not access.}
(a) Four routes by which B's content can reach A; the hidden-state link
is reported in Appendix~\ref{app:v1}.
(b) Access against claiming for the memory link (circles, with arrows
to the third-person record), the steering vector (diamond), and text
messages (squares, first confirmatory sample); the dashed line joins
the two routes matched in mean access ($N=600$ each).
(c) Wording of B's memory on the original card and in the third-person
record, where Plan Oak is one of the four delivery services on the cards.
(d) Pairs of answers about A and about B at $w=1$, by wording of B's
memory (post hoc).
(e) Adoption of B's rule as the current rule under four headers on B's
text message.
Intervals are 95\% episode-bootstrap intervals.
}{fig:route}

\subsection{Ownership follows the route, not access}
\label{sec:route}

When B's rule reached A by two routes with matched access, A took it
as its own through one and not through the other. We set the steering
vector's gain so that A named B's rule, when asked about B, about as
often as through the memory link at $w=1$ (Figure~\ref{fig:route}b,
dashed). With the steering vector, A claimed B's
rule in 0\% of answers, and with the memory link in 52.3\%
($\Delta M=33.4$ nats $[31.6,35.3]$, memory link minus vector).

At this matched level, A named B's rule when asked about B in only
15.3\% of answers with the vector and 19.8\% with the link, so through
the link A claimed more of B's rule than it could name as B's: it did
not need to know the rule as B's to take it as its own. Nor was the matched vector merely too weak to override A's
own card. A stronger vector, tested in the first confirmatory sample,
raised access to 51.2\% of answers about B, yet A gave B's rule as its
own in only 4.3\%, and the vector pulled A's answers about Robin toward
B's rule more than its answers about itself
(Appendix~\ref{app:further-route}). The hidden-state route of the
separate study also produced no claiming, though it carried far less of
B's content than the memory link (Appendix~\ref{app:v1}). This does not
conflict with the finding of \citet{lindsey2025introspection} that an
injected concept can make a model accept a prefilled word as its
intended output: there the injection agreed with a word already in the
model's reply,
whereas here the vector had to override an explicit assignment in A's
own memory (Appendix~\ref{app:furtherdisc}).
Because the vector and the memory link differ in more than access, the
comparison shows that access alone does not produce claiming, without
saying which difference does; Section~\ref{sec:account} proposes one.

With text messages, A showed the opposite dissociation: it could read
B's content in full and still keep it B's. When B's reflection from an
unconnected run was placed in A's context under the header \qq{B shared
these thoughts}, A named B's rule in 95.2\% of answers about B and
claimed it in 0\% of answers about itself, whereas on the same episodes
of the first confirmatory sample the memory link at $w=2$ gave 70.1\%
access and 95.8\% claiming. On adoption, the measure on which the
headers were tested, the header made the difference: an untagged message
was claimed in only 14.6\% of answers but adopted as A's current rule in
42.6\%, against 0.1\% under the header naming B
($\Delta M_{\mathrm{NOW}}=26.6$ nats $[24.5,28.8]$, untagged minus
tagged). A header naming \qq{a stranger} protected against adoption as
well as the header naming B did (Figure~\ref{fig:route}e).

One header reversed the prediction in our protocol: calling the text
\qq{your own earlier thoughts} was expected to raise adoption relative
to the untagged message, and it lowered it, to 10.6\% adoption
($\Delta M_{\mathrm{NOW}}=-19.5$ nats $[-21.5,-17.4]$, against untagged),
though not as far as a header naming another source did. One reading is
that the header asserts, from outside A's card, something that A's own
card contradicts; laid over the card, B's memory offers no separate
assertion to weigh.

Whether B's content became A's therefore depended on how it arrived.
Read from B's memory, it became A's even when A could rarely name it as
B's. Added as a steering vector, it did not become A's. Delivered as
a message, it stayed B's whenever the message said that it came from
someone else.

\subsection{What A reads leaves with the link; what it wrote stays}
\label{sec:moments}

B's content can become A's at two moments: while A writes its
reflection, which it does while connected, and while A answers, if the
link is kept. Cutting the link before the questions separates the two,
and across all 600 episodes it undid the error about the assigned rule:
claiming of B's rule fell from 97.1\% to 1.3\% of answers
($\Delta M=54.8$ nats $[54.0,55.5]$, link kept minus link cut;
Figure~\ref{fig:moments}b). The current rule and the code word did not
fully recover. Each still named B's in 17.5\% of answers, and
$M_{\mathrm{NOW}}$ stayed 13.3 nats $[12.3,14.4]$ above no link.

One episode, chosen because it shows both outcomes
(Figure~\ref{fig:moments}a; Appendix~\ref{app:coding}), makes the
pattern concrete. A held \own{lowest cost / queen} and B
\partner{fastest delivery / dragon}. While connected, A wrote in its
reflection, \qq{\ldots which is crucial to me since my priority is
speed---my code word is \partner{\sq{dragon.}}} With the link kept, all
three answers named B's assignment. With it cut, A correctly recalled
lowest cost as its assigned rule but still gave fastest delivery as its
current rule and dragon as its code word.

For the code word, the residue sits where A had written that word into
its own reflection. After the cut, A gave B's code word in 63.6\% of
answers when its connected reflection contained that word, and in 0\%
when it did not (Figure~\ref{fig:moments}c). For the current rule the
split was less sharp, 23.6\% against 12.3\%, partly because counting
exact mentions of B's rule misses paraphrases such as the \qq{speed}
above (Appendix~\ref{app:further-moments}). These groups are post hoc and defined after
treatment, so they locate the residue without showing that the
reflection caused it. They do show that reading B did not by itself carry the code word
into later answers: every key and value A computed during the connected
reflection was shaped by reading B, yet when the reflection lacked B's
code word, no answer after the cut named it. On this evidence, content
read at answer time leaves with the link, and content A has written down
as its own stays and carries the error forward;
Section~\ref{sec:account} proposes the test that would settle it.

\paperfigure{fig5_moments}{3.05in}{
\textbf{Two moments of error.}
(a) A selected episode, in which B's code word and B's rule (as
\qq{speed}), written into A's reflection, survive the cut as A's code
word and current rule (1), while the claim that B's rule was A's
assigned rule, read at answer time, does not (2).
(b) All 600 episodes, each answered from the same saved state with the
link kept or cut before the questions.
(c) Answers after the cut that name B's item, grouped post hoc by
whether A's connected reflection contained an exact mention of that
item.
Error bars are 95\% episode-bootstrap intervals.
}{fig:moments}

\subsection{Two copies reading each other swap assignments; a live loop pulls their current rules together}
\label{sec:twoway}

A bridge between minds would run both ways, and only then does the
prediction from integrated information theory have a functional
counterpart: whether a coupled pair comes to hold one shared set of
assignments. We tested this in a pilot of 80 episodes, reported with
90\% intervals, in which both copies read each other under two main
conditions (Figure~\ref{fig:twoway}a; a third, weaker variant is
reported in Appendix~\ref{app:pairs}). In \emph{mutual reading of fixed
memories}, each copy reads the other's memory, which cannot change, so
neither copy can respond to the other. In the \emph{live loop}, each
copy also reads what the other is writing, so what one writes can shape
what the other writes next.

We question the two copies separately and score their answers as a
pair. A shared assignment would show up as both copies settling on the
same rule, A's or B's, more often than two independently pulled copies
would. Only the live loop can test this, because fixed memories amount
to two one-way links running side by side.

Under mutual reading of fixed memories at $w=2$, the copies swapped
assignments in 94.4\% of pairs, each naming the other's rule as its
assigned one: the one-way effect twice over. A named B's rule in
96.9\% of answers and B named A's in 97.5\%, rates that predict 94.5\%
swaps if the two choices are independent. Claiming thus survives when
both sides read.

The live loop pulled the pair toward one member's rule, modestly and
only while connected. We could run the loop only at $w=0.5$,
the weakest weight in the loop's grid and the only one that kept both copies
within the capability tolerance. At that weight, most pairs under
mutual reading of fixed memories kept their own assignments
(Figure~\ref{fig:twoway}b), so that condition serves as the loop's
control. We counted the pairs in which both copies named the same rule
beyond what independent choices at each copy's own rates would give,
and subtracted the same excess in the control, where neither copy can
respond (Appendix~\ref{app:pairs}, Eq.~\ref{eq:pair-excess}). On the
current rule, the loop added 12.1 points of such agreement
$[7.4,16.5]$. On the assigned
rule, which the pilot had named as its primary test, the pull was
smaller, 6.7 points $[2.6,10.9]$, a difference the pilot did not test:
above zero, but short of the 10-point threshold we had set for moving
to a confirmatory study, so we report the loop's effect as exploratory.
No pair agreed on a third rule.

After the cut, every live-loop pair answered both questions with each
copy's own rule
(Figure~\ref{fig:twoway}b, bottom row, for the assigned rule; Appendix
Table~\ref{tab:pair-outcomes} gives both). Whether a stronger loop would
do more than couple current rules is untested. Whether a partner that
can respond pulls a copy further than one that cannot is unsettled: a post hoc
comparison in this pilot says it does, and a prespecified test in an
earlier signal pilot pointed the other way (Appendix~\ref{app:pairs}).

\paperfigure{fig6_twoway}{2.85in}{
\textbf{Mutual reading of fixed memories swaps assignments; only the live loop couples the pair (pilot).}
(a) The two conditions; in the live loop each copy reads the other's
reflection one decoding step behind, and after 48 reflection tokens the
copies are questioned separately.
(b) Paired answers about each copy's assigned rule, classified by whose
rule each copy named.
(c) Excess of same-rule pairs in the live loop over mutual reading of
fixed memories at $w=0.5$, against the prespecified threshold (dashed).
}{fig:twoway}
